% Source: physical PDF pp. 420--424 (printed pp. 397--401), from the
% Section 32.3 heading through the final sentence immediately before
% the chapter appendix on physical p. 424.
% Inventory: Eqs. (32.3.1)--(32.3.9), four unnumbered display groups
% (one inside the source footnote), one source footnote converted to an
% ordinary automatically numbered footnote, and linked citation markers
% [4], [8]--[13].
% Corrected two cross-reference ranges in the duality and d=11 discussions.

\subsection{\texorpdfstring{$p$}{p}-Branes}

In some theories in addition to particles there are stable extended
objects, either of infinite extent or stabilized by `wrapping' around a
topologically non-trivial spacetime. The study of supersymmetry and
supergravity in higher-dimensional theories of this sort has opened up
remarkable opportunities for the construction of string theories and
supersymmetric field theories in lower-dimensional spacetime and for the
proof of equivalencies among these
theories,~\hyperref[ch32-ref-4]{[4]},\hyperref[ch32-ref-8]{[8]} which are
beyond the scope of this book. The feature of these extended objects that
concerns us here is that they can carry conserved bosonic quantities other
than those allowed by the Coleman--Mandula theorem. These new conserved
quantities may appear on the right-hand side of the anticommutation
relations of supersymmetry, along with the momentum operator and ordinary
conserved quantities.~\hyperref[ch32-ref-9]{[9]}

In the cases that have been studied so far, the new conserved bosonic
quantities are all \emph{forms}---antisymmetric tensors. For instance, an
object of spatial dimensionality \(p\) (known as a `\(p\)-brane') in a
spacetime of dimensionality \(d\) is described by specifying the \(d\)
spacetime coordinates \(x^\mu(\sigma,t)\) (generally in overlapping
patches covering the object) as functions of the time \(t\) and of a set
of \(p\) coordinates \(\sigma^r\) that parameterize positions on this
object. If the manifold \(x^\mu=x^\mu(\sigma,t)\) at a given time is
topologically non-trivial, in the sense that it cannot be continuously
deformed to a single point, then it may have a non-vanishing value for the
topologically invariant integral\footnote{To see that this integral is
topologically invariant, note that the effect of an infinitesimal change
\(\delta x^\mu(\sigma,t)\) in the function \(x^\mu(\sigma,t)\) is to
change \(I^{\mu_1\mu_2\cdots\mu_p}\) by the amount
\[
  \begin{aligned}
  \delta I^{\mu_1\mu_2\cdots\mu_p}
  ={}&\sum_{n=1}^p\sum_{r_1=1}^p\sum_{r_2=1}^p\cdots\sum_{r_p=1}^p
    \int d\sigma^1\,d\sigma^2\cdots d\sigma^p\,
    \frac{\partial}{\partial\sigma^{r_n}}
    \Biggl[
      \epsilon^{r_1r_2\cdots r_p}\\
  &\quad{}\cdot
      \frac{\partial x^{\mu_1}}{\partial\sigma^{r_1}}
      \frac{\partial x^{\mu_2}}{\partial\sigma^{r_2}}
      \cdots
      \frac{\partial x^{\mu_{n-1}}}{\partial\sigma^{r_{n-1}}}
      \delta x^{\mu_n}
      \frac{\partial x^{\mu_{n+1}}}{\partial\sigma^{r_{n+1}}}
      \cdots
      \frac{\partial x^{\mu_p}}{\partial\sigma^{r_p}}
    \Biggr],
  \end{aligned}
\]
which vanishes when the integral is taken over a compact manifold. It also
vanishes if the integral is over all \(\sigma\), provided
\(\delta x^\mu(\sigma,t)\) is constrained to vanish rapidly when
\(\sigma^r\longrightarrow\infty\).}
\begin{multline}
  I^{\mu_1\mu_2\cdots\mu_p}
  =\int d\sigma^1\,d\sigma^2\cdots d\sigma^p
    \sum_{r_1=1}^p\sum_{r_2=1}^p\cdots\sum_{r_p=1}^p
    \epsilon^{r_1r_2\cdots r_p}\\
  {}\cdot
    \frac{\partial x^{\mu_1}(\sigma,t)}{\partial\sigma^{r_1}}
    \frac{\partial x^{\mu_2}(\sigma,t)}{\partial\sigma^{r_2}}
    \cdots
    \frac{\partial x^{\mu_p}(\sigma,t)}{\partial\sigma^{r_p}}\,.
  \tag{32.3.1}\label{eq:32.3.1}
\end{multline}
The invariance of such integrals under small changes in the functions
\(x^\mu(\sigma,t)\) shows in particular that they are invariant under
spacetime translations, and hence may appear along with \(P^\mu\) and
central charges on the right-hand side of the anticommutation relations
for the supersymmetry generators.~\hyperref[ch32-ref-10]{[10]} The
calculation of the coefficients of such tensors on the right-hand side of
the anticommutation relations is analogous to the Olive--Witten
calculation of the scalar central charge \(Z_{rs}\) in \(N=2\)
supersymmetry theories in four spacetime dimensions, discussed in Section
27.9. We will make no attempt in this section to evaluate these
coefficients or to survey the other non-topological \(p\)-forms that may
appear in the anticommutation relations,~\hyperref[ch32-ref-11]{[11]} but
will simply consider the effects on the supersymmetry algebra of including
conserved antisymmetric tensors that commute with the momentum operators.

It is important that this possibility does not affect the key result that
the supersymmetry generators always belong to the fundamental spinor
representations of the Lorentz group. This is because a totally
antisymmetric tensor in Euclidean coordinates can have at most one
spacetime index equal to 1 and at most one spacetime index equal to \(d\),
and therefore its `weight' defined by Eq. \eqref{eq:32.1.2} can only be
\(\pm1\) or 0. Just as before, this means that the weight of a
supersymmetry generator can only be \(\pm1/2\); Lorentz invariance then
implies that all the \(\sigma\)s defined by Eq. \eqref{eq:32.1.1} are
\(\pm1/2\), which is only possible if the supersymmetry generators belong
to a fundamental spinor representation of \(O(d-1,1)\). Also, because the
new terms in the anticommutators of supersymmetry generators commute with
momentum, the same argument that was given in Section 32.1 shows again
that the supersymmetry generators also commute with momentum.

Lorentz invariance tells us that for non-zero values of the \(p\)-form
`charges,' the anticommutation relations \eqref{eq:32.1.13} and
\eqref{eq:32.1.21}--\eqref{eq:32.1.22} can only take the forms (in the
same notation as in Section 32.1):

\paragraph{\(d\) odd}
\begin{equation}
  \{Q_r,Q_s^{\mathsf T}\}
  =g_{rs}\gamma^\lambda\mathcal C P_\lambda
   +\sum_p z_{rs}^{\mu_1\mu_2\cdots\mu_p}
    \gamma_{\mu_1}\gamma_{\mu_2}\cdots\gamma_{\mu_p}\mathcal C\,.
  \tag{32.3.2}\label{eq:32.3.2}
\end{equation}

\paragraph{\(d\) even}
\begin{multline}
  \{Q_r^\pm,Q_s^{\mp(-1)^{d/2}\,\mathrm T}\}
  =\left(\frac{1\pm\gamma_{d+1}}{2}\right)\\
  {}\cdot
  \left[
    g_{rs}^\pm\gamma^\lambda\mathcal C P_\lambda
    +\sum_{\text{odd }p}
      z_{rs}^{\mu_1\mu_2\cdots\mu_p\,\pm}
      \gamma_{\mu_1}\gamma_{\mu_2}\cdots\gamma_{\mu_p}\mathcal C
  \right].
  \tag{32.3.3}\label{eq:32.3.3}
\end{multline}
\begin{multline}
  \{Q_r^\pm,Q_s^{\pm(-1)^{d/2}\,\mathrm T}\}
  =\left(\frac{1\pm\gamma_{d+1}}{2}\right)\\
  {}\cdot
    \sum_{\text{even }p}
      z_{rs}^{\mu_1\mu_2\cdots\mu_p\,\pm}
      \gamma_{\mu_1}\gamma_{\mu_2}\cdots\gamma_{\mu_p}\mathcal C\,.
  \tag{32.3.4}\label{eq:32.3.4}
\end{multline}
(Recall that \(\mathcal C\) appears in the anticommutation relations
because
\(\mathscr J_{\mu\nu}^{\mathsf T}
=-\mathcal C^{-1}\mathscr J_{\mu\nu}\mathcal C\); the \(Q_r^\pm\) are
supersymmetry generators in the case of even \(d\) for which
\(\gamma_{d+1}Q_r^\pm=\pm Q_r^\pm\); and
\(\mathcal C^{-1}\gamma_{d+1}\mathcal C
=(-1)^{d/2}\gamma_{d+1}\).) For even \(d\) we have
\[
  \epsilon^{\mu_1\mu_2\cdots\mu_d}
  \gamma_{\mu_1}\gamma_{\mu_2}\cdots\gamma_{\mu_p}
  \mathrel{\propto}
  \gamma_{d+1}\gamma_{\mu_{p+1}}\gamma_{\mu_{p+2}}
  \cdots\gamma_{\mu_d}\,,
\]
while for \(d\) odd
\[
  \epsilon^{\mu_1\mu_2\cdots\mu_d}
  \gamma_{\mu_1}\gamma_{\mu_2}\cdots\gamma_{\mu_p}
  \mathrel{\propto}
  \gamma_{\mu_{p+1}}\gamma_{\mu_{p+2}}\cdots\gamma_{\mu_d}\,.
\]
Thus \(p\)-branes and \(d-p\) branes make the same contributions in
Eqs. \eqref{eq:32.3.2}--\eqref{eq:32.3.4} for any \(d\), so that we can restrict \(p\) to
run only over values from 0 to \(d/2\) for \(d\) even and from 0 to
\((d-1)/2\) for \(d\) odd.

The symmetry of anticommutators is reflected in symmetry conditions on
the \(p\)-brane central charges \(z_{rs}^p\) in
Eqs. \eqref{eq:32.3.2}--\eqref{eq:32.3.4}. Eqs. \eqref{eq:32.A.15} and
\eqref{eq:32.A.30} give
\begin{equation}
  \gamma_\mu^{\mathsf T}
  =(-1)^n\mathcal C^{-1}\gamma_\mu\mathcal C\,,
  \qquad
  \mathcal C^{\mathsf T}
  =(-1)^{n(n+1)/2}\mathcal C\,,
  \tag{32.3.5}\label{eq:32.3.5}
\end{equation}
for both \(d=2n\) and \(d=2n+1\). These give the antisymmetrized products
\(\gamma_{[\mu_1}\gamma_{\mu_2}\cdots\gamma_{\mu_p]}\) the symmetry
property
\begin{align}
  \gamma_{[\mu_1}\gamma_{\mu_2}\cdots\gamma_{\mu_p]}\mathcal C
  &=
  (-1)^{pn}(-1)^{n(n+1)/2}
  \left[
    \gamma_{[\mu_p}\gamma_{\mu_{p-1}}\cdots\gamma_{\mu_1]}\mathcal C
  \right]^{\mathsf T}
  \notag\\
  &=
  (-1)^{pn}(-1)^{n(n+1)/2}(-1)^{p(p-1)/2}
  \left[
    \gamma_{[\mu_1}\gamma_{\mu_2}\cdots\gamma_{\mu_p]}\mathcal C
  \right]^{\mathsf T}.
  \tag{32.3.6}\label{eq:32.3.6}
\end{align}
It follows immediately that for odd \(d\),
\begin{equation}
  z_{rs}^{\mu_1\mu_2\cdots\mu_p}
  =(-1)^{pn}(-1)^{n(n+1)/2}(-1)^{p(p-1)/2}
   z_{sr}^{\mu_1\mu_2\cdots\mu_p}\,,
  \tag{32.3.7}\label{eq:32.3.7}
\end{equation}
while for even \(d\),
\begin{equation}
  z_{rs}^{\mu_1\mu_2\cdots\mu_p\,\pm}
  =(-1)^{pn}(-1)^{n(n+1)/2}(-1)^{p(p-1)/2}
   z_{sr}^{\mu_1\mu_2\cdots\mu_p\,(-1)^n(-1)^p\mp}\,.
  \tag{32.3.8}\label{eq:32.3.8}
\end{equation}

Consider for instance the important case of \(N=1\) supersymmetry in
\(d=11\) spacetime dimensions, which is one version of the popular
\(M\)-theory generalization of string theories. Eq. \eqref{eq:32.3.7} shows that
the single \(p\)-form central charge \(z^{\mu_1\mu_2\cdots\mu_p}\)
vanishes unless
\begin{equation}
  -(-1)^{pn}(-1)^{p(p-1)/2}=+1\,,
  \tag{32.3.9}\label{eq:32.3.9}
\end{equation}
which is satisfied only for \(p\) equal to 1, 2, and 5. The value \(p=1\)
is realized by the momentum operator itself, which arises from particles
as well as extended objects. The other possibilities, \(p=2\) and
\(p=5\), arise in theories with 2-branes and 5-branes, respectively. Note
that there can be no other independent tensor central charges, such as a
1-form arising from 1-branes, because the number of independent components
in \(P^\mu\) and in a 2-form and a 5-form is
\[
  11+\binom{11}{2}+\binom{11}{5}=528\,,
\]
while the number of independent components in an anticommutator of two
32-component fundamental spinors is \(32\times33/2=528\).

Just as the 0-form electric charge is the source of a 1-form gauge field
\(A_\mu(x)\), so also a \(p\)-form conserved quantity
\(z^{\mu_1\mu_2\cdots\mu_p}\) may serve as the source of a \(p+1\)-form
gauge field \(A_{\mu_1\mu_2\cdots\mu_{p+1}}\) of the sort discussed in
Section 8.8. In fact, such gauge fields do appear in supergravity
theories. For instance, as remarked in the previous section, the \(N=1\)
supergravity theory in \(d=11\) spacetime dimensions includes a massless
particle whose states are in the 3-form representation of the \(O(9)\)
little group, and therefore must be described by a 3-form gauge field
\(A_{\mu\nu\rho}(x)\). The study of solutions of this supergravity theory
shows that there are two-branes~\hyperref[ch32-ref-12]{[12]} that indeed
do provide sources for \(A_{\mu\nu\rho}(x)\). Also, as noted in Section
8.8, this gauge theory is equivalent to one with a
\((d-p-2=6)\)-form gauge field, which can have the 5-form
\(z^{\mu_1\cdots\mu_5}\) as a source, and there are indeed 5-brane
solutions which provide such sources for this 6-form gauge
field.~\hyperref[ch32-ref-13]{[13]} The \(N=1\) eleven-dimensional
supersymmetry algebra does in fact receive contributions from these
2-branes and 5-branes.~\hyperref[ch32-ref-11]{[11]}
